% Generated draft: requires mathematical and source review.
\documentclass[11pt]{article}
\usepackage[margin=25mm]{geometry}
\usepackage{fontspec}
\setmainfont{Latin Modern Roman}
\usepackage{amsmath,amssymb,mathtools,unicode-math}
\setmathfont{Latin Modern Math}
\usepackage{enumitem,needspace,xurl,hyperref}
\setlength{\parindent}{0pt}
\setlength{\parskip}{5pt}
\setlength{\emergencystretch}{3em}
\newcommand{\N}{\mathbb{N}}
\newcommand{\Z}{\mathbb{Z}}
\newcommand{\Q}{\mathbb{Q}}
\newcommand{\R}{\mathbb{R}}
\newcommand{\C}{\mathbb{C}}
\newcommand{\F}{\mathbb{F}}
\newcommand{\A}{\mathbb{A}}
\newcommand{\PP}{\mathbb{P}}
\newcommand{\E}{\mathbb{E}}
\newcommand{\eps}{\varepsilon}
\newcommand{\dd}{\,\mathrm{d}}
\begin{document}
\section*{UnsolvedMath ID: 30000318}
\textbf{Source problem code:} catalogue entry
\subsection{Definitions and mathematical statement}
\textbf{Question.} Must every maximal abelian self-adjoint subalgebra \(A\) of the free group factor \(M=L(\F_2)\) satisfy \(\infty\in\operatorname{Puk}(A\subset M)\)? Here \(\F_2\) is the free group on two generators, and \(L(\F_2)\) is the von Neumann algebra generated by its left regular representation on \(\ell^2(\F_2)\), with canonical trace \(\tau\). A masa means a unital abelian von Neumann subalgebra with \(A'\cap M=A\). The quantifier is over all such masas, not only a constructed or otherwise distinguished subclass.\par In the standard definition of the Puk\'anszky invariant, work on \(\mathcal H=L^2(M,\tau)\), let \(J\) be the canonical conjugation \(J(x)=x^*\) on the dense copy of \(M\), and let \(e_A\) project onto \(L^2(A,\tau)\). Put \(\mathcal B=(A\cup JAJ)''\), generated by the left and right actions of \(A\). The type-I algebra \((1-e_A)\mathcal B'(1-e_A)\) has homogeneous type-\(\mathrm I_n\) summands, with \(n\in\N_{>0}\cup\{\infty\}\). The set of \(n\) for which the corresponding summand is nonzero is \(\operatorname{Puk}(A\subset M)\). The requested conclusion is that its type-\(\mathrm I_\infty\) part is nonzero. It is stronger than unboundedness of the finite multiplicities or the existence of infinitely many distinct finite entries. The corner removes the canonical \(L^2(A)\) component.\par The supplied literature triage checked 2026-08-21 labels the question open, reporting that the source explicitly records it and that later examples do not settle the universal statement. No precise later example or theorem is reproduced, so the current status must be independently checked. The definition above clarifies the standard invariant without claiming progress on the all-masas question.
\subsection{Short English statement}
Does every masa in the two-generator free group factor have a genuinely infinite-multiplicity component in its Puk\'anszky invariant? The claim concerns all masas, and an actual \(\infty\) entry is different from arbitrarily large finite entries.
\subsection{Sources}
\begin{itemize}\item Local UnsolvedMath registry, permanent ID 30000318, code catalogue entry; \emph{Puk\'anszky Invariants of Masas in \(L(F_2)\)}, Oberwolfach Reports Open Problems.\item \emph{C*-Algebren (2006)}, Oberwolfach Reports, DOI 10.4171/owr/2005/41, pp. 2305--2374. \url{https://doi.org/10.4171/owr/2005/41}. Proposed year and year tag: 2005; no proposer named.\item Supplied literature label \emph{original Oberwolfach report}: \url{https://doi.org/10.4171/owr/2005/41}.\item Supplied literature label \emph{Pukansky invariants of masas}: \url{https://www.researchgate.net/publication/245573920_C-Algebren}. No separate bibliographic theorem citation is supplied with this label.\item Literature triage checked 2026-08-21, status \texttt{open}; record updated 2026-09-16. The \texttt{corrected\_verified} field concerns restoration of free-group and infinity notation. Verify dated claims independently.\end{itemize}
\end{document}
